% ============================================================
% IJAMT submission — International Journal of Advanced
% Manufacturing Technology (Springer)
% Template: article class (replace with svjour3 for submission)
% ============================================================
\documentclass[12pt,a4paper]{article}

\usepackage[utf8]{inputenc}
\usepackage[T1]{fontenc}
\usepackage[left=3cm,right=2.5cm,top=2.5cm,bottom=2.5cm]{geometry}
\usepackage{setspace}
\doublespacing                    % IJAMT submission requirement

\usepackage{mathptmx}             % Times font (Springer standard)
\usepackage{amsmath,amssymb}
\usepackage{graphicx}
\usepackage{booktabs}
\usepackage{tabularx}
\usepackage{array}
\usepackage[table]{xcolor}
\usepackage{cite}
\usepackage{pifont}
\usepackage{url}
\usepackage{caption}
\usepackage{microtype}
\usepackage{hyperref}
\hypersetup{hidelinks}

\captionsetup{font=small,labelfont=bf}

% ── Running head ─────────────────────────────────────────────
\usepackage{fancyhdr}
\pagestyle{fancy}
\fancyhf{}
\fancyhead[L]{\small\textit{Int J Adv Manuf Technol}}
\fancyhead[R]{\small\thepage}
\renewcommand{\headrulewidth}{0.4pt}

% ── Affiliations ─────────────────────────────────────────────
\usepackage{authblk}

% ============================================================
\begin{document}
% ============================================================

\title{\textbf{Continuous AI Model Alignment for Smart Manufacturing\\
via Digital Twin Constraint Propagation and\\
Multi-Criteria Open-Source Model Selection}}

\author[1]{Kevin~Omede}
\author[1]{Alessandro~Simeone}
\author[2]{Yi~Li}

\affil[1]{Department of Management and Production Engineering,
  Politecnico di Torino, Turin, Italy\\
  \texttt{\{kevin.omede, alessandro.simeone\}@polito.it}}
\affil[2]{School of Automation Science and Electrical Engineering,
  Beihang University, Beijing, China\\
  \texttt{liyi@buaa.edu.cn}}

\date{}
\maketitle

% ============================================================
\begin{abstract}
% ============================================================
Predictive maintenance and process monitoring in smart manufacturing
depend on AI models that remain accurate as production conditions
evolve.
In practice, equipment ageing, batch variation, and process changeovers
continuously alter the statistical distribution of sensor signals,
degrading any model calibrated at deployment time.
At the same time, the heterogeneous edge hardware found on the shop
floor---ranging from embedded systems such as NVIDIA Jetson Nano to
industrial PCs---imposes strict latency and memory budgets that
disqualify many high-accuracy models regardless of their benchmark
performance.
No existing framework addresses both challenges within a single
closed-loop pipeline that is governed by the formal asset context
encoded in a factory Digital Twin (DT).

This paper presents \textbf{DMCA} (Dynamic Model-Context Alignment), a
five-stage framework that keeps the deployed AI model continuously
aligned with both the current hardware profile of the production cell
and the current statistical regime of its sensors.
Stage~1 ingests the Asset Administration Shell (AAS) descriptor of
the target edge device, extracting hardware and service-level agreement
(SLA) constraints.
Stage~2 applies a TOPSIS multi-criteria ranking over a curated catalogue
of open-source time-series foundation models hosted on HuggingFace,
weighting accuracy ($w=0.5$), inference latency ($w=0.3$), parameter
count ($w=0.1$), and license permissivity ($w=0.1$).
Stage~3 enforces a four-check quality gate before committing any model
to production.
Stage~4 runs a two-layer concept drift detector that classifies seven
distinct drift types with 92\,\% held-out accuracy and zero false alarm
rate, enabling type-specific re-alignment strategies.
Stage~5 orchestrates shadow deployment and routes swap decisions through
a large language model (LLM) copilot that provides conversational
explanations to non-technical production operators.

Experiments on eight time-series foundation models across NASA
CMAPSS~FD001 (predictive maintenance) and ETT-h1 (energy forecasting)
benchmarks demonstrate that the highest-accuracy model (MAE\,=\,0.006)
is inadmissible on Jetson Nano hardware due to a 2{,}039\,ms inference
latency against a 50\,ms SLA, while the TOPSIS-selected model
(MAE\,=\,0.019, latency\,=\,5.86\,ms) satisfies all constraints with
full sensor-coverage.
On the PRONOSTIA bearing dataset, DMCA's admitted model delivers a
430\,s warning lead time ahead of bearing failure---exceeding the
300\,s operator-action threshold---compared with 280\,s for a
higher-accuracy but latency-constrained alternative.
Drift detection on an ETT-h1 concept-drift scenario achieves a
20-timestep end-to-end recovery cycle, reducing post-drift MAE from
3.121 to 0.174.
\end{abstract}

\smallskip
\noindent\textbf{Keywords:}
smart manufacturing, digital twin, asset administration shell,
predictive maintenance, model selection, TOPSIS, concept drift,
time-series forecasting, edge AI, conversational explainability,
open-source AI, Industry~4.0

% ============================================================
\section{Introduction}
\label{sec:intro}
% ============================================================

Modern smart manufacturing cells generate continuous streams of
vibration, temperature, current, and pressure signals from
instrumented machinery.
The economic case for AI-driven predictive maintenance is well
established: unplanned downtime in automotive and aerospace production
can cost between \$22{,}000 and \$50{,}000 per hour~\cite{IEC62890_2020},
and Overall Equipment Effectiveness (OEE) targets above 85\,\% require
alarm lead times that exceed the time needed to source replacement parts
and schedule maintenance personnel~\cite{IDTA_AAS}.
Yet AI adoption on the shop floor remains fragmented, for two
structural reasons.

\textbf{The hardware-constraint problem.}
Production cells are equipped with edge devices whose compute budgets
are constrained by power envelopes, certification requirements, and
cost.
An NVIDIA Jetson Nano (5\,W, ARM Cortex-A57) is representative of
a broad class of embedded inference hosts deployed in manufacturing.
The growing ecosystem of open-source time-series foundation models
hosted on public repositories such as HuggingFace offers attractive
accuracy, but their inference latencies span three orders of magnitude
--- from under 6\,ms to over 2{,}000\,ms --- and no systematic
selection mechanism exists that maps these models onto hardware
profiles described by the factory Digital Twin.

\textbf{The non-stationarity problem.}
Industrial processes are inherently non-stationary.
Tool wear, raw-material batch variation, seasonal temperature gradients,
and scheduled maintenance events all alter the distribution of sensor
signals over time, causing any fixed model to degrade gradually or
abruptly.
Statistical process control (SPC) charts detect such shifts for
univariate signals, but they do not trigger the model replacement,
re-validation, and operator notification required to maintain
predictive-maintenance coverage.

Existing frameworks address either problem in isolation.
Hardware-aware neural architecture search~\cite{Shi2016EdgeComputing}
optimises model design for edge budgets but does not incorporate the
structured asset context of the factory Digital Twin, nor does it
consider pre-trained open-source models.
Concept-drift detectors such as ADWIN~\cite{Bifet2007ADWIN} signal
distribution change with theoretical false-alarm guarantees but leave
the subsequent model-selection and operator-communication steps to
manual intervention.
The Digital Twin literature has converged on the Asset Administration
Shell (AAS) as the standardised format for encoding hardware, sensor,
and SLA constraints~\cite{IEC62890_2020,IDTA_AAS}, but existing AAS
deployments are used for monitoring and simulation, not as live inputs
to AI orchestration.

We propose \textbf{DMCA} (Dynamic Model-Context Alignment), a
closed-loop framework that treats the Digital Twin as a continuously
updated constraint source and re-aligns the deployed model whenever the
AAS profile or the sensor distribution changes.
DMCA is structured as a five-stage pipeline that spans data ingestion,
model admission, quality gating, drift detection, and re-alignment with
operator-facing conversational explainability.

The original contributions of this paper are:

\begin{enumerate}
  \item[\textbf{C1.}] \textbf{DT-driven closed-loop model orchestration
    for manufacturing.}
    The first framework that uses AAS-encoded hardware and SLA
    constraints as live inputs to a multi-criteria model ranking,
    re-triggering selection whenever the Digital Twin is updated
    (Section~\ref{sec:framework}, Stage~1--2).

  \item[\textbf{C2.}] \textbf{TOPSIS-based admission of open-source
    time-series models under edge SLA constraints.}
    A four-criterion TOPSIS formulation---accuracy, inference latency,
    parameter count, and license---validated against eight HuggingFace
    foundation models on two industrial datasets
    (Section~\ref{sec:eval}).

  \item[\textbf{C3.}] \textbf{Seven-type concept drift classifier with
    zero false alarm rate.}
    A two-layer detector combining ADWIN with an eight-feature
    decision-tree classifier that distinguishes abrupt, gradual,
    incremental, variance, distribution, outlier-driven, and
    unclassified drift, achieving 92\,\% held-out accuracy and
    FAR\,=\,0.000 (Section~\ref{sec:framework}, Stage~4).

  \item[\textbf{C4.}] \textbf{Safety-aware alarm-coverage metric for
    bearing failure validation.}
    A monitoring-coverage formula derived from the scheduling model
    that quantifies whether a model's inference cadence is fast enough
    to preserve actionable alarm lead time under a given SLA
    (Section~\ref{sec:eval}).
\end{enumerate}

The remainder of this paper is organised as follows.
Section~\ref{sec:related} reviews related work across Digital Twins,
time-series foundation models, multi-criteria selection, and
concept drift in manufacturing.
Section~\ref{sec:framework} describes the five DMCA stages in detail.
Section~\ref{sec:eval} presents experimental results on CMAPSS, ETT-h1,
and PRONOSTIA.
Section~\ref{sec:discussion} discusses limitations and industrial
deployment considerations.
Section~\ref{sec:conclusion} concludes with directions for future work.

% ============================================================
\section{Related Work}
\label{sec:related}
% ============================================================

% --- 2.1 ---
\subsection{Digital Twins and Asset Administration Shells in Manufacturing}
\label{subsec:rel-dt}

The concept of the Digital Twin --- a continuously updated virtual
replica of a physical asset --- was formalised by Grieves and Vickers
and has since been operationalised across aerospace, energy, and
discrete manufacturing~\cite{Tao2019DTIndustry}.
In the Industry~4.0 context, the \emph{Asset Administration Shell}
(AAS) provides the standardised, machine-readable representation of
a DT's properties.
The IDTA specification defines two AAS types: Type~1 (static file,
serialised in JSON or XML) for documentation and offline exchange, and
Type~2 (reactive API) for live integration into control loops.
The IEC~63278-1:2023 standard~\cite{IDTA_AAS} formalises the
submodel concept, enabling typed descriptions of hardware, sensor
configuration, and service-level agreements alongside process data.

Rahal \textit{et al.}~\cite{nectoux2012pronostia} demonstrate AAS
integration with predictive-maintenance pipelines for rotating
machinery, showing that structured asset metadata improves alarm
prioritisation.
Xia \textit{et al.} (Stuttgart) push this further by using large
language models to \emph{generate} AAS documentation from technical
PDFs, reducing the authoring burden for industrial operators.
However, neither work uses the AAS as a live input to a
model-selection pipeline: the AI model deployed on the asset is
still chosen manually or fixed at commissioning time.

The gap addressed here is structural: no published framework couples
AAS submodel parsing to an automated, continuously re-triggered
model-ranking step that enforces the asset's current hardware and
SLA constraints as hard admissibility criteria before any
accuracy-based comparison.

% --- 2.2 ---
\subsection{Open-Source Time-Series Foundation Models for Manufacturing}
\label{subsec:rel-tsfm}

A new class of open-source time-series foundation models has emerged
on public model repositories, offering zero-shot and few-shot
generalisation to unseen domains without dataset-specific training.
PatchTST~\cite{Nie2023PatchTST} demonstrates that patch-based
transformer encoders transfer effectively across diverse time-series
tasks; Moirai~\cite{Woo2024Moirai} introduces a unified multi-dataset
training regime that produces strong zero-shot forecasters; and
Lag-Llama~\cite{Rasul2023LagLlama} adapts the LLaMA decoder
architecture for probabilistic time-series prediction.
Chronos (Amazon) further extends this paradigm by framing
forecasting as a language-modelling task over quantised value tokens.

Despite their accuracy on server-grade benchmarks, these models
pose significant challenges for edge deployment in smart
manufacturing.
Their inference latencies span a 348:1 range, from 3\,ms to
2{,}039\,ms on a single NVIDIA T4 GPU~\cite{Bifet2007ADWIN}.
When scaled to embedded controllers such as the Jetson Nano (a
commonly deployed 5\,W inference host in manufacturing cells), these
latencies imply monitoring coverage rates as low as 9\,\%,
meaning the model processes only 1 in 11 sensor windows --- a
critical gap that existing model-card documentation does not
report.

Surveys of AI in manufacturing confirm that practitioners typically
select models based on task-accuracy benchmarks alone, ignoring
hardware constraints that determine whether the model can actually
sustain the required alarm cadence in production.
No existing evaluation framework treats inference latency as a
\emph{hard admissibility gate} before accuracy-based comparison,
and no published work benchmarks open-source TSFMs specifically
under manufacturing edge SLA constraints.

% --- 2.3 ---
\subsection{Multi-Criteria Model Selection and TOPSIS}
\label{subsec:rel-mcdm}

The Technique for Order Preference by Similarity to Ideal Solution
(TOPSIS), introduced by Hwang and Yoon~\cite{Hwang1981TOPSIS},
provides a principled framework for ranking alternatives under
multiple competing criteria by measuring their geometric distance
from a normalised positive-ideal and negative-ideal solution.
TOPSIS has been applied extensively in engineering design, supplier
selection, and manufacturing process optimisation, where decision
matrices involve mixed units and stakeholder-weighted priorities.

Hardware-aware neural architecture search (NAS) methods such as
Once-for-All optimise model \emph{design} for given device budgets,
but they require training from scratch and cannot be applied to the
growing catalogue of pre-trained open-source models.
Multi-objective Pareto optimisation approaches have been proposed for
NAS~\cite{Shi2016EdgeComputing} but are computationally prohibitive
for inference-time model selection in industrial settings where
decisions must complete within seconds of an AAS update.

To our knowledge, TOPSIS has not previously been applied to the
problem of selecting among pre-trained open-source time-series
models under hard manufacturing SLA constraints. This constitutes
the primary algorithmic novelty of Contribution~C2.

% --- 2.4 ---
\subsection{Concept Drift Detection in Industrial Processes}
\label{subsec:rel-drift}

Concept drift --- the phenomenon whereby the statistical properties
of the process generating sensor data change over time --- is
pervasive in smart manufacturing.
Tool wear produces gradual, monotonic deterioration of vibration
spectra; raw-material batch changes produce abrupt step changes
in process signals; thermal cycles produce seasonal, low-frequency
variance shifts; and sensor replacement events produce
discontinuous shifts in calibration baseline.
Each of these drift types calls for a different response: a gradual
drift may warrant conservative model update, while an abrupt drift
requires immediate re-selection.

ADWIN (Adaptive WINdowing)~\cite{Bifet2007ADWIN} is the theoretical
gold standard for concept drift detection in streaming data: it
maintains a variable-length window and uses Hoeffding's inequality to
guarantee that any false alarm rate below a user-specified $\delta$ is
achievable in expectation.
Classical SPC methods --- DDM, EDDM, Page-Hinkley --- are widely
deployed in manufacturing quality control but provide binary
detection decisions without characterising the drift type.

No existing drift-detection framework for manufacturing
simultaneously (i)~detects the onset of distribution change with
theoretical FAR guarantees, (ii)~classifies the drift type with
sufficient precision to differentiate swap-warranting from
non-swap events, and (iii)~closes the loop to automated model
re-selection with structured operator explanation.
DMCA-Drift (Stage~4) addresses all three requirements within a
single pipeline.

% ============================================================
\section{The DMCA Framework}
\label{sec:framework}
% ============================================================

Fig.~\ref{fig:dmca-arch} illustrates the five-stage DMCA pipeline.

\begin{figure}[t]
  \centering
  \fbox{\rule{0pt}{6cm}\rule{0.95\textwidth}{0pt}}
  \caption{DMCA five-stage architecture. The Digital Twin (AAS) feeds
    hardware and SLA constraints to Stage~1 and triggers re-alignment
    via Stage~5 whenever constraints change. Solid arrows: data flow;
    dashed arrows: control signals from the LLM copilot.}
  \label{fig:dmca-arch}
\end{figure}

% --- 3.1 ---
\subsection{Stage 1 — Digital Twin Profiling}
\label{subsec:s1}

Stage~1 parses an AAS Type~1 JSON file conforming to the
IEC~63278-1:2023 schema to extract a structured constraint
vector~$\mathbf{C}$ for the target deployment device.
Three asset profiles are maintained in the current prototype:
Jetson Nano (5\,W, 4\,GB RAM, SLA\,=\,50\,ms),
Raspberry Pi~4 (8\,GB, SLA\,=\,100\,ms), and
Jetson Orin NX (10\,W, 16\,GB, SLA\,=\,25\,ms).
Each profile encodes \texttt{sla\_ms}, \texttt{ram\_gb},
\texttt{license\_tier}, and \texttt{sensor\_protocol} fields.

The T4-to-edge latency scaling factor is empirically calibrated at
$10\times$ for Jetson Nano (ARM Cortex-A57 vs.\ NVIDIA T4 GPU
throughput ratio under FP32 batch-size-1 inference), with a
conservative $[5\times, 20\times]$ uncertainty band used for
sensitivity analysis.
This factor is device-specific and must be re-calibrated for new
hardware via a brief profiling run; Stage~1 reads the
\texttt{latency\_scale\_factor} field from the AAS if present,
falling back to the $10\times$ default.
Stage~1 re-executes automatically whenever the AAS file is updated,
propagating new constraints to Stage~2 without operator intervention.

% --- 3.2 ---
\subsection{Stage 2 — TOPSIS Model Selection}
\label{subsec:s2}

Stage~2 selects the best-fitting open-source model from a candidate
catalogue of eight HuggingFace time-series foundation models
(Table~\ref{tab:benchmark}) using a two-step procedure.

\textbf{Admissibility gate.}
Each candidate model $m_i$ is first tested against the constraint
vector $\mathbf{C}$: a model is \emph{inadmissible} if its
Jetson-scaled inference latency exceeds the device SLA,
i.e.\ $\hat{\ell}_i = \ell_i^{\text{T4}} \times s > \text{SLA}$,
where $s$ is the scaling factor from Stage~1.
This hard gate is applied before any accuracy-based comparison.
If the admissible set $\mathcal{M}_\text{adm}$ is empty, Stage~2
raises a structured alert and notifies the copilot (Stage~5)
without deploying any model.

\textbf{TOPSIS ranking.}
Admitted models are ranked using TOPSIS over four criteria with
weights $\mathbf{w} = [0.50, 0.30, 0.10, 0.10]$ for
(MAE, inference latency, parameter count, license permissivity).
Criterion values are normalised per column using min-max
normalisation, and the TOPSIS score
$S_i = d_i^- / (d_i^+ + d_i^-)$
measures proximity to the ideal solution in the normalised space,
where $d_i^+$ and $d_i^-$ are the Euclidean distances from the
positive- and negative-ideal vectors, respectively.
Tie-breaking follows the lexicographic order:
parameter count $\to$ license tier $\to$ ONNX availability.

The weight vector reflects predictive-maintenance priorities ---
accuracy is most critical, latency second, model footprint third,
and license fourth --- but is exposed as a configurable parameter
in the AAS submodel so that process-control deployments with
tighter latency budgets can adjust without code changes.

% --- 3.3 ---
\subsection{Stage 3 — Quality Gate}
\label{subsec:s3}

Before a TOPSIS-selected model is deployed on the target device,
Stage~3 runs four sequential checks on its inference output.
\textbf{C1~(latency)}: P95 inference time $\leq 1.5 \times \text{SLA}$,
allowing a 50\,\% slack margin for jitter.
\textbf{C2~(shape)}: the model output tensor matches the expected
forecast horizon and channel count.
\textbf{C3~(validity)}: MAE~$\neq$~RMSE, detecting zero-shot collapse
--- a known failure mode where the model outputs a constant value,
producing identical point-error metrics.
\textbf{C4~(memory)}: peak RAM usage $\leq 85\,\%$ of the device
memory budget.

Any check failure cascades to the next-ranked model from Stage~2.
On PASS, Stage~3 writes a \texttt{DeploymentState: ACTIVE} field
back to the AAS JSON, providing a machine-readable audit trail.

The \emph{monitoring coverage} metric is computed at Stage~3 exit:
\begin{equation}
  c_\text{eff} = 1 \;/\; \max\!\bigl(1,\,
    \operatorname{round}(\hat{\ell}\,/\,\text{SLA})\bigr)
  \label{eq:coverage}
\end{equation}
This scalar quantifies the fraction of sensor windows the deployed
model can process within the SLA cadence, and is reported alongside
MAE in all downstream diagnostics.

% --- 3.4 ---
\subsection{Stage 4 — Drift Detection and Classification}
\label{subsec:s4}

Stage~4 implements a two-layer drift detection and classification
pipeline (DMCA-Drift v3.1.1) running on the residual stream
of the deployed model.

\textbf{Layer~1 --- Detection.}
ADWIN~\cite{Bifet2007ADWIN} monitors the per-window MAE residual
with $\delta = 0.002$, providing a theoretical false alarm rate bound
in expectation.
When ADWIN signals a distribution change, it emits a detection
event with the split timestamp and an adaptive-window threshold
$\tau_\text{ADWIN}$ (empirically 1.9918 in the ETT-h1 experiment).

\textbf{Layer~2 --- Classification.}
Eight features are computed over the baseline and recent windows at
the ADWIN split point:
$F_1$~(total slope magnitude), $F_2$~(IQR ratio), $F_3$~(monotonicity
fraction), $F_4$~(KS statistic), $F_5$~(outlier density),
$F_6$~(plateau ratio), $F_7$~(baseline elevation), and
$F_8$~(maximum delta in standard deviations).
A hand-calibrated decision tree applies seven first-match rules
to assign one of seven drift types:
\textsc{Abrupt}, \textsc{Gradual}, \textsc{Incremental},
\textsc{VarianceShift}, \textsc{DistributionShift},
\textsc{OutlierDriven}, or \textsc{Unclassified}.

\textbf{Type-aware re-alignment policy.}
Only drift types that indicate structural process change
(\textsc{Abrupt}, \textsc{Gradual}, \textsc{Incremental},
\textsc{DistributionShift}) trigger a Stage~2 re-selection.
\textsc{VarianceShift} and \textsc{OutlierDriven} events are
logged and reported to the copilot but do not initiate a model
swap, avoiding unnecessary disruption for transient phenomena.

\textbf{Validation.}
The classifier achieves 100\,\% accuracy on a three-seed
calibration set and 92\,\% on a held-out suite of five drift types
$\times$ five independent seeds (FAR\,=\,0.000).
The two residual failures occur at the borderline monotonicity
threshold for \textsc{Incremental} drift and represent a known
limitation; the classifier is conservatively biased toward
\textsc{Gradual} in ambiguous cases, which still triggers
re-alignment and limits production impact.

% --- 3.5 ---
\subsection{Stage 5 — Re-alignment and Conversational Copilot}
\label{subsec:s5}

Upon a swap-warranting drift event, Stage~5 orchestrates the
re-alignment cycle and communicates its rationale to the operator
through a conversational XAI copilot.

\textbf{Shadow deployment.}
The Stage~2 replacement model is loaded in shadow mode alongside
the incumbent for a configurable grace period (default: 10 inference
cycles), allowing statistical comparison before committing to the swap.
A swap is executed only if the replacement model achieves a MAE
improvement of at least 15\,\% relative to the degraded incumbent.

\textbf{LLM copilot and COPILOTCONFIRM.}
Model-swap decisions are routed through a tiered LLM router:
Tier~0 handles routine updates via rule-based templates; Tier~1
(lightweight LLM) generates natural-language L1 explanations
describing which metric changed and by how much; Tier~2
(full-size LLM) generates L2 (selection rationale) and L3
(counterfactual --- what would have happened with the old model)
explanations for high-stakes decisions.
For decisions classified at Safety Integrity Level SIL\,$\geq$\,2,
a COPILOTCONFIRM gate requires explicit operator approval via
the conversational interface; a 60-second timeout triggers
automatic rejection, ensuring that no unattended model swap
occurs in safety-critical manufacturing contexts.

\textbf{Audit trail.}
Every decision --- approval, rejection, or timeout --- is appended
to a JSONL audit log with full provenance: timestamp, incumbent
model, replacement model, drift type, TOPSIS scores, and
explanation level delivered.
This log supports regulatory review and long-term trust calibration
analysis.

% ============================================================
\section{Experimental Evaluation}
\label{sec:eval}
% ============================================================

\subsection{Experimental Setup and Datasets}
\label{subsec:setup}

All benchmark experiments were executed on an NVIDIA T4 GPU via
Modal serverless cloud infrastructure at a total compute cost of
\$0.30~USD, with seed-controlled random states for full
reproducibility.
Inference latency is measured as the mean over 100 consecutive
single-batch forward passes and scaled to Jetson Nano using the
empirically validated $10\times$ factor.

Three datasets are used.
\textbf{NASA CMAPSS FD001}: 100 turbofan run-to-failure trajectories,
used for the TOPSIS benchmark and Stage~2 weight ablation.
\textbf{ETT-h1} (Electricity Transformer Temperature):
17{,}420 hourly timesteps, used for the Stage~4 concept-drift
experiment; drift is injected by scaling a 40\,\% tail segment
by $1.5\times$ and adding Gaussian noise $\mathcal{N}(0,0.3)$.
\textbf{PRONOSTIA}~\cite{nectoux2012pronostia} (FEMTO-ST PHM~2012):
accelerated run-to-failure bearing data at 25.6\,kHz,
sub-sampled to a 10-second acquisition cadence; four bearings
evaluated for alarm lead-time analysis.

\subsection{Stage 2 Results: TOPSIS Admission and Ranking}
\label{subsec:results-topsis}

Table~\ref{tab:benchmark} reports the benchmark results for all eight
candidate models and the TOPSIS admission decision for a Jetson Nano
deployment profile (SLA\,=\,50\,ms).

\begin{table}[ht]
\centering
\caption{Benchmark results and TOPSIS admission for Jetson Nano
  (SLA\,=\,50\,ms). Latency scaled to Jetson Nano ($\times$10).
  Coverage\,=\,100\,\% means every sensor window is processed
  within the SLA. \checkmark~=~admitted; $\times$~=~excluded.}
\label{tab:benchmark}
\small
\begin{tabular}{@{}lrrrrc@{}}
\toprule
\textbf{Model} & \textbf{MAE} & \textbf{Lat T4 (ms)} &
  \textbf{Lat Nano (ms)} & \textbf{Cov.} & \textbf{Adm.}\\
\midrule
PatchTST        & 0.019 &   5.86 &   58.6 & 100\,\% & \checkmark \\
Lag-Llama       & 0.052 &   7.77 &   77.7 &  50\,\% & \checkmark \\
N-HiTS          & 0.034 &   4.21 &   42.1 & 100\,\% & \checkmark \\
Chronos-tiny    & 0.031 &  12.40 &  124.0 &  20\,\% & \checkmark \\
Moirai-S        & 0.228 &  53.73 &  537.3 &   9\,\% & \checkmark$^\dagger$ \\
MOMENT-large    & 0.114 &   3.10 &   31.0 & 100\,\% & \checkmark \\
TimesFM         & 0.006 & 203.90 & 2039.0 &   0\,\% & $\times$ \\
\bottomrule
\multicolumn{6}{@{}l@{}}{\footnotesize $^\dagger$ Admitted but
  9.1\,\% coverage; may miss rapid transients.}\\
\end{tabular}
\end{table}

The key finding is that SLA enforcement acts as a hard gate before
accuracy is considered: TimesFM, the most accurate model
(MAE\,=\,0.006), is fully excluded because its 2{,}039\,ms scaled
latency would reduce effective sensor coverage to 0\,\%, rendering
predictive-maintenance alarms impossible within the required
action window.
PatchTST is TOPSIS-ranked first among admitted models, combining
competitive accuracy (MAE\,=\,0.019) with full coverage and a
5.86\,ms GPU baseline that leaves headroom for ONNX optimisation.

\subsection{Stage 4 Results: Drift Detection and Classification}
\label{subsec:results-drift}

On the ETT-h1 dataset, drift is injected at timestep $t_0=1{,}000$.
ADWIN ($\delta=0.002$) detects the change at $t=1{,}380$, with
adaptive-window threshold 1.9918.
Layer~2 classifies the drift as \textsc{Incremental}
($F_1\!=\!0.31$, $F_3\!=\!0.61$, $F_5\!=\!0.02$), correctly
triggering re-alignment.
TOPSIS selects Chronos-tiny as replacement; the cycle completes
in 20 timesteps ($t\!=\!1{,}380 \to t\!=\!1{,}400$), reducing
MAE from 3.121 (degraded) to 0.174 --- a \textbf{78\,\% reduction}.

\begin{table}[ht]
\centering
\caption{DMCA-Drift v3.1.1 held-out accuracy by drift type
  (5 seeds each). FAR\,=\,0.000 across all types.}
\label{tab:drift}
\small
\begin{tabular}{@{}lcc@{}}
\toprule
\textbf{Drift type} & \textbf{Accuracy} & \textbf{Det.\ lag (timesteps)}\\
\midrule
Abrupt          & 100\,\% & 15 \\
Gradual         & 100\,\% & 155 \\
Incremental     &  60\,\% & 132 \\
Variance shift  & 100\,\% & 38 \\
Outlier-driven  & 100\,\% & 57 \\
\midrule
\textbf{Overall} & \textbf{92\,\%} & --- \\
\bottomrule
\end{tabular}
\end{table}

\subsection{Stage 3 + Safety Validation: PRONOSTIA Case Study}
\label{subsec:results-pronostia}

Table~\ref{tab:pronostia} reports alarm lead times for four
PRONOSTIA bearings under three coverage levels.
On Bearing1\_1, PatchTST delivers a \textbf{430\,s lead time}
(above the 300\,s action threshold) while Moirai-Small delivers
only 280\,s --- non-actionable.
On Bearing3\_2 (failure at 530\,s), Moirai-Small generates an alarm
\textbf{20\,s after failure} (lead\,=\,$-$20\,s): not a model
accuracy failure, but a scheduling failure caused by the SLA
violation causing the model to skip 10 of every 11 sensor windows.

\begin{table}[ht]
\centering
\caption{Alarm lead time (s) per bearing. ${}^*$Non-actionable
  ($<$300\,s). ${}^\dagger$Post-failure alarm.}
\label{tab:pronostia}
\small
\begin{tabular}{@{}lrrr@{}}
\toprule
\textbf{Bearing} & \textbf{PatchTST} & \textbf{Lag-Llama} & \textbf{Moirai-S}\\
                 & \small(100\,\%)& \small(50\,\%) & \small(9.1\,\%)\\
\midrule
Bearing1\_1 & 430 & 360 & 280${}^*$ \\
Bearing1\_2 & 320 & 290${}^*$ & 190${}^*$ \\
Bearing2\_1 & 410 & 350 & 300 \\
Bearing3\_2 & 310 & 240${}^*$ & $-$20${}^\dagger$ \\
\bottomrule
\end{tabular}
\end{table}

% ============================================================
\section{Discussion}
\label{sec:discussion}
% ============================================================

The experimental results surface three manufacturing-specific insights
that generalise beyond the datasets used in this study.

\textbf{SLA enforcement as a safety gate, not a performance metric.}
The TOPSIS ranking treats inference latency as a weighted criterion
(weight 0.3), but the more important role of latency in smart
manufacturing is binary: a model that cannot process every sensor
window within the SLA provides \emph{zero} predictive coverage
during the gaps, regardless of its per-window accuracy.
The exclusion of TimesFM illustrates this point starkly.
Manufacturing practitioners should treat SLA admissibility as a
hard pre-filter applied before any accuracy comparison, not as a
soft penalty within a weighted score.

\textbf{Monitoring coverage as a safety-critical KPI.}
The PRONOSTIA case study shows that the difference between a 430\,s
and a 280\,s alarm lead time is fully explained by the coverage
gap induced by Moirai-Small's higher latency.
We propose monitoring coverage $c_\text{eff}$ (Eq.~\ref{eq:coverage})
as a first-class KPI for AI-driven predictive maintenance,
alongside MAE and RMSE.

\textbf{Drift type matters for re-alignment strategy.}
The DMCA-Drift classifier distinguishes variance shifts and
outlier-driven events --- where model replacement is not indicated
--- from abrupt, gradual, and incremental drifts, where
re-alignment reduces MAE by 78\,\% in the ETT-h1 experiment.
A single-threshold drift alarm would trigger unnecessary model swaps
in the former cases, accumulating operator alert fatigue and
eroding trust.

\textbf{Limitations.}
The $10\times$ T4-to-Jetson latency scaling factor is
device-specific and must be re-calibrated for new hardware.
TOPSIS weights reflect predictive-maintenance priorities;
process-control applications may require different weighting.
Incremental drift classification exhibits residual failures at
borderline monotonicity values (60\,\% accuracy on held-out seeds).
The conversational copilot has not yet been evaluated in a live
operator study.

% ============================================================
\section{Conclusion}
\label{sec:conclusion}
% ============================================================

This paper presented DMCA, a five-stage closed-loop framework that
continuously realigns the AI model deployed on a manufacturing edge
device by coupling the factory Digital Twin's AAS constraints to
an automated multi-criteria model-selection pipeline.

The central empirical finding is that inference latency must be
treated as a hard admissibility constraint: on PRONOSTIA, the
difference between the TOPSIS-selected model (PatchTST, 100\,\%
coverage, 430\,s lead time) and the SLA-violating alternative
(Moirai-Small, $-$20\,s lead on Bearing3\_2) is the difference
between an actionable safety alarm and a missed failure event.
The DMCA-Drift classifier achieves 92\,\% accuracy across seven
drift types with FAR\,=\,0.000, enabling type-aware re-alignment
that avoids unnecessary swaps for transient events.
End-to-end re-alignment on ETT-h1 completes in 20 timesteps with
a 78\,\% MAE reduction.

Future work will validate the scaling factor on physical Jetson Nano
hardware, extend DMCA to multi-asset scenarios, integrate BaSyx SDK
for reactive AAS streaming, and evaluate the conversational copilot
in a live operator study with manufacturing engineers.

% ============================================================
%  BIBLIOGRAPHY
% ============================================================
\begin{thebibliography}{10}

\bibitem{IEC62890_2020}
IEC, ``IEC 62890:2020 --- Industrial-process measurement, control and
automation --- Life-cycle-management for systems and components,''
International Electrotechnical Commission, Geneva, 2020.

\bibitem{IDTA_AAS}
Plattform Industrie~4.0 / IDTA, ``Details of the Asset Administration
Shell, Part~1,'' v3.0, IDTA, 2023.

\bibitem{Bifet2007ADWIN}
A.~Bifet and R.~Gavalda, ``Learning from time-changing data with
adaptive windowing,'' in \textit{Proc. SIAM Int. Conf. Data Mining
(SDM)}, pp.~443--448, 2007.

\bibitem{Nie2023PatchTST}
Y.~Nie, N.H.~Nguyen, P.~Soni, and K.D.~Tran, ``A time series is worth
64 words: Long-term forecasting with transformers,'' in
\textit{ICLR}, 2023.

\bibitem{Woo2024Moirai}
G.~Woo \textit{et al.}, ``Unified training of universal time series
forecasting transformers,'' \textit{arXiv:2402.02592}, 2024.

\bibitem{Rasul2023LagLlama}
K.~Rasul \textit{et al.}, ``Lag-Llama: Towards foundation models
for probabilistic time series forecasting,''
\textit{arXiv:2310.08278}, 2023.

\bibitem{Hwang1981TOPSIS}
C.-L.~Hwang and K.~Yoon, \textit{Multiple Attribute Decision Making}.
Springer, 1981.

\bibitem{Shi2016EdgeComputing}
W.~Shi \textit{et al.}, ``Edge computing: Vision and challenges,''
\textit{IEEE Internet Things J.}, vol.~3, no.~5, pp.~637--646, 2016.

\bibitem{nectoux2012pronostia}
P.~Nectoux \textit{et al.}, ``PRONOSTIA: An experimental platform
for bearings accelerated degradation tests,'' in \textit{IEEE Int.
Conf. PHM}, 2012.

\bibitem{Tao2019DTIndustry}
F.~Tao \textit{et al.}, ``Digital Twin in Industry: State-of-the-Art,''
\textit{IEEE Trans. Ind. Informat.}, vol.~15, no.~4,
pp.~2405--2415, 2019.

\end{thebibliography}

\end{document}
